\section{Basic properties}

\begin{proposition}[Monotone smoothness reparameterization]
\label{prop:monotone}
For positive penalized eigenvalues $\delta_1,\ldots,\delta_{L-d}$,
\begin{equation}
\label{eq:smoothness_first_derivative}
S'(\lambda)
=
\frac{1}{L-d}
\sum_{j=1}^{L-d}
\frac{\delta_j}{(1+\lambda\delta_j)^2}
>0,
\end{equation}
and
\begin{equation}
\label{eq:smoothness_second_derivative}
S''(\lambda)
=
-\frac{2}{L-d}
\sum_{j=1}^{L-d}
\frac{\delta_j^2}{(1+\lambda\delta_j)^3}
<0.
\end{equation}
After adjoining $\lambda=\infty$, $S$ maps $[0,\infty]$ continuously and strictly increasingly onto $[0,1]$.
\end{proposition}

\begin{proof}
Differentiate \cref{eq:S} term by term. Positivity of each $\delta_j$ gives the signs; the endpoint values follow from $(1+\lambda\delta_j)^{-1}$.
\end{proof}

The bounded scale has a common interpretation across fixed $(d,L)$ configurations while preserving the ordering of penalty strength.

\begin{proposition}[Effective degrees of freedom]
\label{prop:edf}
For fixed $\lambda$,
\begin{equation}
\label{eq:edf_smoothness_formula}
\operatorname{tr}(H_\lambda)=L-(L-d)S(\lambda).
\end{equation}
\end{proposition}

\begin{proof}
From \cref{eq:H_spectral},
\begin{equation}
\label{eq:edf_spectral_trace}
\operatorname{tr}(H_\lambda)
=
d+\sum_{j=1}^{L-d}(1+\lambda\delta_j)^{-1}.
\end{equation}
Substitution of \cref{eq:S} yields the result.
\end{proof}

\begin{proposition}[Equivalent optimization in penalty and smoothness space]
\label{prop:equiv}
Let $F(S)=f(\lambda(S))$. Then the global minimizers in $\lambda\in[0,\infty]$ and $S\in[0,1]$ correspond under the bijection in Proposition \cref{prop:monotone}. At interior points,
\begin{equation}
\label{eq:loss_change_of_variables}
F'(S)=\frac{f'(\lambda)}{S'(\lambda)},
\end{equation}
so interior stationary points correspond exactly.
\end{proposition}

\begin{proof}
A strictly increasing bijection preserves the minimizing set. The derivative identity follows from the chain rule.
\end{proof}

\begin{proposition}[Exact endpoints and the unique limiting fit]
\label{prop:endpoints}
Under the conditions of \cref{prop:pls_wellposed}, $H(0)=I$.
Let $U_0\in\mathbb R^{L\times d}$ have orthonormal columns spanning
$\ker(D_d)$. At $S=1$, interpreted as $\lambda\to\infty$,
\begin{equation}
\label{eq:exact_limiting_projection}
H(1)=H_\infty=U_0U_0^\top=P_{\ker(D_d)}.
\end{equation}
Although $H_\infty$ is singular with rank $d$ and nullity $L-d$,
the limiting trend is uniquely defined by
\begin{equation}
\label{eq:exact_limiting_constrained_fit}
H_\infty x_T
=\operatorname*{arg\,min}_{\tau:\,D_d\tau=0}
\|x_T-\tau\|_2^2.
\end{equation}
Thus no $1-\varepsilon$ surrogate is required for $S=1$.
\end{proposition}

\begin{proof}
At $\lambda=0$, $H_0=I$. By \cref{eq:H_spectral}, the
$d$ unit eigenvalues on $\ker(D_d)$ remain one, whereas
$(1+\lambda\delta_j)^{-1}\to0$ for every $\delta_j>0$.
This proves the projection formula. Any feasible
$\tau\in\ker(D_d)$ has the form $\tau=U_0a$.
Since $U_0^\top U_0=I$, the least-squares objective
$\|x_T-U_0a\|_2^2$ has the unique minimizer
$a=U_0^\top x_T$. Hence $\tau=U_0U_0^\top x_T$ is unique,
even though the projection matrix itself is not invertible.
\end{proof}

For $d=2$, the endpoints are the unsmoothed historical path and
its unique least-squares linear projection. The expression
$(I+\infty Q)^{-1}$ is not an ordinary matrix inverse:
$S=1$ must be computed as the exact spectral projection.
A large finite penalty or $S=1-\varepsilon$ is only an
approximation and can introduce avoidable numerical error.
Endpoint comparison is part of the scientific selection problem,
not merely a numerical detail.

\subsection{The future polynomial family}

The estimator produces a polynomial continuation only after the
future-difference rule has been specified. This matters for
interpreting forecast-CV as a prospective trend-selection method.

\begin{proposition}[Forecast polynomial indexed by smoothness]
\label{prop:continuation_polynomial}
Fix $1\le d<L$ and $h\ge1$. For every $S\in[0,1]$, the
continuation $k\mapsto\widehat\tau_{T+k\mid T}(S)$ is the
restriction to positive integer steps of a polynomial in
$k$ of degree at most $d-1$. The coefficients depend on
$S$ through the terminal backward differences of $H(S)x_T$.
\end{proposition}

\begin{proof}
For a fitted sequence $u=H(S)x_T$, imposing zero future
$d$th differences gives the discrete Newton continuation
\begin{equation}
\label{eq:newton_forecast_polynomial}
\widehat\tau_{T+k\mid T}(S)
=\sum_{j=0}^{d-1}\binom{k+j-1}{j}
(\nabla^j u)_L.
\end{equation}
Each binomial coefficient is a polynomial in $k$ of
degree $j$. The sum therefore has degree at most
$d-1$. Since $u=H(S)x_T$, its terminal differences
and consequently its coefficients vary with $S$.
The limiting case $S=1$ is well defined by the
orthogonal projection in \cref{prop:endpoints}.
\end{proof}

This proposition does not imply that the degree $d-1$
is chosen by forecast-CV: the baseline fixes $d$ and
selects among the corresponding extrapolations by
their historical future-block error. Changing $d$
would require an explicitly specified joint or
nested model-selection experiment.

\subsection{Existence and possible nonuniqueness}

The fitting problem has a unique solution for every $S\in[0,1]$: for $S<1$ this follows from strict convexity in \cref{prop:pls_wellposed}, and for $S=1$ from the constrained least-squares argument in \cref{prop:endpoints}. For finite $S<1$, $H_{\lambda(S)}$ depends continuously on $S$; the exact projection limit at $S=1$ makes the full family continuous on $[0,1]$. Thus, with a fixed finite collection of historical folds, $F^{\mathrm{pool}}_{T,d,L,h}(S)$ is continuous on the compact interval $[0,1]$ and attains at least one global minimum. \emph{Existence does not imply a unique minimizing $S$.}

Uniqueness is not assumed. Distinct smoothing regimes can produce similar validation loss, and the criterion may contain several local minima. This is statistically meaningful because competing minima can imply different effective degrees of freedom and terminal slopes. The implemented numerical search uses adaptive derivative bracketing, Brent root refinement, and exact endpoint comparison. The mathematical criterion in \cref{eq:Sstar} is independent of numerical search details and does not require tracking minima through time.

\subsection{Forecast optimality versus recovery optimality}

If a latent trend $\tau$ is known in simulation, a recovery-oriented oracle might choose
\begin{equation}
\label{eq:latent_recovery_oracle}
S_{\mathrm{rec}}
\in
\arg\min_S
\|\widehat\tau_S-\tau\|^2.
\end{equation}
The proposed criterion instead chooses $S^\star_{d,L,h}$ from future predictive error. These targets need not coincide. A smoother can reconstruct the historical latent path accurately but extrapolate poorly, or sacrifice historical recovery in exchange for more stable terminal derivatives.

This distinction is a central empirical question of the paper, not a claim that forecast-optimal smoothness is universally preferable under every loss.
